\section{Diện tích hình phẳng}

\subsection{Nền tảng lý thuyết \& Định lý cơ bản}

\begin{hopdinhly}[Định lý]
Cho hàm số $y=f(x)$ liên tục trên đoạn $[a;b]$. Diện tích $S$ của hình phẳng giới hạn bởi đồ thị hàm số $y=f(x)$, trục hoành $Ox$ và hai đường thẳng $x=a$, $x=b$ được tính bởi công thức
\[
S = \int_a^b \ababs{f(x)} \dx.
\]
\end{hopdinhly}

\begin{center}
\hinhdientich{-0.5*(\x-1)*(\x-3)+1}{1}{3}{4}{3.2}
\end{center}

\begin{hopdinhly}[Định lý]
Cho hai hàm số $y=f(x)$, $y=g(x)$ liên tục trên đoạn $[a;b]$. Diện tích $S$ của hình phẳng giới hạn bởi hai đồ thị $y=f(x)$, $y=g(x)$ và hai đường thẳng $x=a$, $x=b$ được tính bởi công thức
\[
S = \int_a^b \ababs{f(x)-g(x)} \dx.
\]
\end{hopdinhly}

\begin{center}
\hinhhaihdc{-0.35*(\x-1)^2+3}{0.4*(\x-1)^2+0.3}{0}{2.5}{3.2}{4}
\end{center}

\begin{hopchuy}[Bản chất của dấu giá trị tuyệt đối]
Tích phân $\displaystyle\int_a^b f(x)\dx$ có thể mang giá trị âm nếu $f(x)<0$ trên $[a;b]$, nhưng \textbf{diện tích luôn không âm}. Đây chính là lý do công thức diện tích luôn xuất hiện dấu trị tuyệt đối -- và cũng là nguồn gốc của toàn bộ kỹ thuật xử lý trong 4 dạng toán dưới đây.
\end{hopchuy}

\newpage
\subsection{Dạng 1: Tính diện tích khi giả thiết cho sẵn hàm số và cận}

\begin{hopphuongphap}
\begin{enumerate}[leftmargin=1cm]
  \item Giải phương trình $f(x)=0$ (hoặc $f(x)=g(x)$) để tìm các nghiệm thuộc $(a;b)$.
  \item Các nghiệm này chia $[a;b]$ thành các đoạn con mà trên mỗi đoạn, $f(x)$ (hoặc $f(x)-g(x)$) không đổi dấu.
  \item Dùng một giá trị đại diện trong mỗi đoạn con để xét dấu, từ đó phá dấu trị tuyệt đối.
  \item Tính từng tích phân thành phần rồi cộng lại.
\end{enumerate}
\end{hopphuongphap}

\begin{hopvidu}[(Hàm đa thức bậc hai)]
Tính diện tích hình phẳng giới hạn bởi parabol $(P): y=x^2-4x+3$, trục hoành và hai đường thẳng $x=0$, $x=3$.
\end{hopvidu}
\begin{loigiai}
Xét phương trình hoành độ giao điểm với $Ox$: $x^2-4x+3=0 \Leftrightarrow x=1$ hoặc $x=3$. Nghiệm $x=1$ thuộc $(0;3)$ nên ta xét dấu trên hai đoạn $[0;1]$ và $[1;3]$.

Với $x\in[0;1]$: chọn $x=0,5$, ta có $f(0{,}5)=0{,}25-2+3=1{,}25>0$.
Với $x\in[1;3]$: chọn $x=2$, ta có $f(2)=4-8+3=-1<0$.

Vậy
\[
S=\int_0^1 (x^2-4x+3)\dx + \int_1^3 -(x^2-4x+3)\dx.
\]
Ta có $\displaystyle\int (x^2-4x+3)\dx = \frac{x^3}{3}-2x^2+3x+C$. Do đó
\[
\int_0^1(x^2-4x+3)\dx = \left.\left(\frac{x^3}{3}-2x^2+3x\right)\right|_0^1 = \frac{4}{3},
\qquad
\int_1^3(x^2-4x+3)\dx = 0-\frac{4}{3}=-\frac{4}{3}.
\]
\ketqua{S=\frac{4}{3}+\frac{4}{3}=\frac{8}{3} \ (\text{đvdt})}
\end{loigiai}

\begin{hopvidu}[(Hàm đa thức bậc ba)]
Tính diện tích hình phẳng giới hạn bởi đồ thị $y=x^3-3x$, trục hoành và hai đường thẳng $x=-1$, $x=2$.
\end{hopvidu}
\begin{loigiai}
Xét $x^3-3x=0 \Leftrightarrow x(x-\sqrt3)(x+\sqrt3)=0 \Leftrightarrow x\in\{-\sqrt3;0;\sqrt3\}$. Trong đoạn $[-1;2]$ chỉ có hai nghiệm $x=0$ và $x=\sqrt3$ thỏa mãn, chia $[-1;2]$ thành ba đoạn con.

Xét dấu: trên $[-1;0]$, tại $x=-0{,}5$: $f=-0{,}5(0{,}25-3)=1{,}375>0$; trên $[0;\sqrt3]$, tại $x=1$: $f=1-3=-2<0$; trên $[\sqrt3;2]$, tại $x=1{,}9$: $f\approx 1{,}16>0$.

Gọi $F(x)=\dfrac{x^4}{4}-\dfrac{3x^2}{2}$ là một nguyên hàm của $f(x)=x^3-3x$. Ta tính:
\[
\int_{-1}^{0} f(x)\dx = F(0)-F(-1)=0-\left(-\frac{5}{4}\right)=\frac{5}{4},
\]
\[
\int_{0}^{\sqrt3} f(x)\dx = F(\sqrt3)-F(0)=-\frac{9}{4}-0=-\frac{9}{4}
\ \Rightarrow\ \text{lấy } \frac{9}{4},
\]
\[
\int_{\sqrt3}^{2} f(x)\dx = F(2)-F(\sqrt3)=-2-\left(-\frac{9}{4}\right)=\frac{1}{4}.
\]
\ketqua{S=\frac{5}{4}+\frac{9}{4}+\frac{1}{4}=\frac{15}{4} \ (\text{đvdt})}
\end{loigiai}

\begin{hopvidu}[(Hàm phân thức)]
Tính diện tích hình phẳng giới hạn bởi đồ thị $\displaystyle y=\frac{x-2}{x+1}$, trục hoành và hai đường thẳng $x=0$, $x=3$.
\end{hopvidu}
\begin{loigiai}
Ta có $\dfrac{x-2}{x+1}=0 \Leftrightarrow x=2 \in (0;3)$, và hàm số liên tục trên $[0;3]$ (vì $-1\notin[0;3]$). Viết lại
\[
f(x)=\frac{x-2}{x+1}=1-\frac{3}{x+1}, \qquad F(x)=x-3\ln|x+1|.
\]
Xét dấu: tại $x=0$, $f(0)=-2<0$; tại $x=2{,}5$, $f(2{,}5)=\dfrac{0{,}5}{3{,}5}>0$. Vậy $f<0$ trên $[0;2]$ và $f>0$ trên $[2;3]$.
\[
S = -\int_0^2 f(x)\dx + \int_2^3 f(x)\dx = -\big[F(2)-F(0)\big] + \big[F(3)-F(2)\big] = F(0)+F(3)-2F(2).
\]
Thay số: $F(0)=0$, $F(2)=2-3\ln 3$, $F(3)=3-3\ln4$, ta được
\ketqua{S=6\ln3-3\ln4-1 \approx 1{,}43 \ (\text{đvdt})}
\end{loigiai}

\begin{hopvidu}[(Hàm lượng giác)]
Tính diện tích hình phẳng giới hạn bởi đồ thị hàm số $y=\sin x$, trục hoành và hai đường thẳng $x=0$, $x=2\pi$.
\end{hopvidu}
\begin{loigiai}
Trên $[0;2\pi]$, $\sin x=0$ tại $x=0,\pi,2\pi$; ta có $\sin x\ge0$ trên $[0;\pi]$ và $\sin x\le 0$ trên $[\pi;2\pi]$.
\[
S=\int_0^\pi \sin x\dx - \int_\pi^{2\pi}\sin x\dx = \big[-\cos x\big]_0^\pi - \big[-\cos x\big]_\pi^{2\pi} = 2-(-2)=4.
\]
\ketqua{S=4 \ (\text{đvdt})}
\end{loigiai}

\begin{hopchuy}[Sai lầm thường gặp]
Rất nhiều học sinh quên xét dấu và viết thẳng $S=\int_a^b f(x)\dx$ mà không đưa trị tuyệt đối vào, dẫn đến kết quả \textbf{âm hoặc sai lệch} khi đồ thị cắt trục hoành ở giữa đoạn $[a;b]$ (như Ví dụ 2, 3, 4 ở trên). Luôn kiểm tra: \emph{nghiệm của $f(x)=0$ có thuộc khoảng $(a;b)$ hay không?}
\end{hopchuy}

\newpage
\subsection{Dạng 2: Tính diện tích khi thiếu cận (giới hạn bởi các đường cong khép kín)}

\begin{hopphuongphap}
\begin{enumerate}[leftmargin=1cm]
  \item Giải phương trình hoành độ giao điểm $f(x)=g(x)$ để tìm các cận $a,b$ (không có cận cho sẵn trong đề bài).
  \item Nếu có nhiều hơn 2 nghiệm phân biệt trong miền xét, hình phẳng bị chia thành nhiều miền con -- phải \textbf{chẻ tích phân} theo từng đoạn giữa hai nghiệm liên tiếp.
  \item Tính diện tích từng miền rồi cộng lại như Dạng 1.
\end{enumerate}
\end{hopphuongphap}

\begin{hopvidu}[(Parabol và đường thẳng)]
Tính diện tích hình phẳng giới hạn bởi parabol $(P): y=x^2-2x$ và đường thẳng $d: y=x$.
\end{hopvidu}
\begin{loigiai}
Phương trình hoành độ giao điểm: $x^2-2x=x \Leftrightarrow x^2-3x=0 \Leftrightarrow x=0$ hoặc $x=3$. Đây chính là hai cận $a=0,\,b=3$.

Xét dấu $h(x)=x-(x^2-2x)=3x-x^2$ trên $(0;3)$: tại $x=1$, $h(1)=2>0$, nên $d$ nằm trên $(P)$.
\[
S=\int_0^3 (3x-x^2)\dx = \left.\left(\frac{3x^2}{2}-\frac{x^3}{3}\right)\right|_0^3 = \frac{27}{2}-9.
\]
\ketqua{S=\frac{9}{2} \ (\text{đvdt})}
\end{loigiai}

\begin{hopvidu}[(Hai miền tách biệt -- 3 giao điểm)]
Tính diện tích hình phẳng giới hạn bởi đồ thị $y=x^3$ và đường thẳng $y=x$.
\end{hopvidu}
\begin{loigiai}
Phương trình hoành độ giao điểm: $x^3-x=0 \Leftrightarrow x(x-1)(x+1)=0 \Leftrightarrow x\in\{-1;0;1\}$. Ba nghiệm này chia miền thành \textbf{hai} hình phẳng riêng biệt: $[-1;0]$ và $[0;1]$.

Xét dấu $h(x)=x^3-x$: trên $(-1;0)$, tại $x=-0{,}5$: $h=-0{,}125+0{,}5=0{,}375>0$; trên $(0;1)$, tại $x=0{,}5$: $h=0{,}125-0{,}5=-0{,}375<0$.
\[
S=\int_{-1}^{0}(x^3-x)\dx + \int_0^1 \big[-(x^3-x)\big]\dx.
\]
Với $F(x)=\dfrac{x^4}{4}-\dfrac{x^2}{2}$: $\displaystyle\int_{-1}^0(x^3-x)\dx = F(0)-F(-1)=0-\left(-\frac14\right)=\frac14$, và $\displaystyle\int_0^1(x^3-x)\dx=F(1)-F(0)=-\frac14$, nên miền thứ hai đóng góp $\frac14$.
\ketqua{S=\frac14+\frac14=\frac12 \ (\text{đvdt})}
\end{loigiai}

\begin{center}
\begin{tikzpicture}[scale=1, declare function={f(\x)=\x^3; g(\x)=\x;}]
  \begin{axis}[
      axis lines=middle, xlabel={$x$}, ylabel={$y$},
      xmin=-1.6, xmax=1.6, ymin=-1.6, ymax=1.6,
      xtick={-1,0,1}, ytick=\empty, width=8cm, height=8cm,
      axis line style={-Latex}, samples=100]
    \addplot [domain=-1:0, teal, fill=teal, fill opacity=0.3] {f(x)} \closedcycle;
    \addplot [domain=0:1, amber, fill=amber, fill opacity=0.3] {f(x)} \closedcycle;
    \addplot [domain=-1.5:1.5, navy, thick] {f(x)} node[right]{$y=x^3$};
    \addplot [domain=-1.5:1.5, navylight, thick] {g(x)} node[above left]{$y=x$};
  \end{axis}
\end{tikzpicture}
\end{center}

\begin{hopvidu}[(Hai parabol)]
Tính diện tích hình phẳng giới hạn bởi hai parabol $y=x^2$ và $y=2-x^2$.
\end{hopvidu}
\begin{loigiai}
Phương trình hoành độ giao điểm: $x^2=2-x^2 \Leftrightarrow 2x^2=2 \Leftrightarrow x=\pm1$. Tại $x=0$: $x^2=0 < 2-x^2=2$, vậy $y=2-x^2$ nằm trên $y=x^2$ trên $(-1;1)$.
\[
S=\int_{-1}^1 \big[(2-x^2)-x^2\big]\dx = \int_{-1}^1(2-2x^2)\dx = \left.\left(2x-\frac{2x^3}{3}\right)\right|_{-1}^{1}.
\]
\ketqua{S=\frac{4}{3}+\frac{4}{3}=\frac{8}{3} \ (\text{đvdt})}
\end{loigiai}

\begin{hopvidu}[(Công thức tính nhanh: Parabol -- Cát tuyến)]
Tính diện tích hình phẳng giới hạn bởi parabol $(P): y=-x^2+2x+3$ và đường thẳng $d: y=x-1$.
\end{hopvidu}
\begin{loigiai}
Đặt $h(x)=f(x)-g(x)=(-x^2+2x+3)-(x-1)=-x^2+x+4$. Phương trình $h(x)=0$ có $\Delta = 1+16=17>0$ nên có hai nghiệm phân biệt $x_1,x_2$ với
\[
x_2-x_1 = \frac{\sqrt\Delta}{|a|}=\sqrt{17} \quad (a=-1).
\]
\begin{hopchuy}[Công thức nhanh]
Khi $h(x)=a(x-x_1)(x-x_2)$ (parabol và đường thẳng, hoặc hai parabol cùng hệ số $x^2$ trái dấu), diện tích được tính nhanh bởi
\[
S=\frac{|a|}{6}(x_2-x_1)^3,
\]
không cần tính cụ thể $x_1,x_2$.
\end{hopchuy}
Áp dụng: $|a|=1$, do đó
\ketqua{S=\frac{1}{6}\left(\sqrt{17}\right)^3=\frac{17\sqrt{17}}{6} \ (\text{đvdt})}
\end{loigiai}

\newpage
\subsection{Dạng 3: Kỹ năng đọc hiểu đồ thị và tư duy hình học}

\begin{hopphuongphap}
Khi đề bài chỉ cho \textbf{hình vẽ} (không cho biểu thức $f(x)$ tường minh), ta không thể giải phương trình $f(x)=0$ để xét dấu. Thay vào đó:
\begin{enumerate}[leftmargin=1cm]
  \item Quan sát trực tiếp trên hình: phần đồ thị nằm \textbf{trên} hay \textbf{dưới} trục hoành (hoặc đường cong kia) trên từng khoảng.
  \item Dùng dữ kiện $\displaystyle\int$ đã cho (nếu có) để suy ra diện tích từng miền, có để ý dấu.
  \item Cộng các miền diện tích (luôn dương) lại với nhau.
\end{enumerate}
\end{hopphuongphap}

\begin{hopvidu}[(Phân biệt tích phân và diện tích)]
Cho hàm số $y=f(x)$ liên tục trên $[-2;4]$ có đồ thị như hình vẽ. Biết diện tích miền tô màu xanh (phía trên trục $Ox$, trên $[-2;1]$) bằng $3$ và diện tích miền tô màu cam (phía dưới trục $Ox$, trên $[1;4]$) bằng $5$. Tính diện tích hình phẳng giới hạn bởi đồ thị $y=f(x)$, trục $Ox$ và hai đường $x=-2$, $x=4$.
\end{hopvidu}

\begin{center}
\begin{tikzpicture}[scale=0.95]
  \draw[-Latex] (-2.6,0) -- (4.6,0) node[right]{$x$};
  \draw[-Latex] (0,-2.2) -- (0,2.6) node[above]{$y$};
  \draw[navy, thick, domain=-2:4, samples=80, smooth]
     plot (\x, {1.6*sin(deg(\x*0.85)) *exp(-0.05*\x)});
  \fill[teal, opacity=0.3, domain=-2:1, samples=80]
     (-2,0) -- plot[smooth, domain=-2:1] (\x,{1.6*sin(deg(\x*0.85))*exp(-0.05*\x)}) -- (1,0) -- cycle;
  \fill[amber, opacity=0.3, domain=1:4, samples=80]
     (1,0) -- plot[smooth, domain=1:4] (\x,{1.6*sin(deg(\x*0.85))*exp(-0.05*\x)}) -- (4,0) -- cycle;
  \node[teal] at (-0.6,0.7) {$S_1=3$};
  \node[amber] at (2.5,-0.9) {$S_2=5$};
  \draw[dashed, shadowgray] (-2,0)--(-2,-0.3) node[below]{$x=-2$};
  \draw[dashed, shadowgray] (4,0)--(4,-0.3) node[below]{$x=4$};
  \draw[dashed, shadowgray] (1,0)--(1,1.4);
\end{tikzpicture}
\end{center}

\begin{loigiai}
Vì diện tích luôn dương nên ta \textbf{không được} cộng trực tiếp các tích phân có dấu -- ta cộng các diện tích thành phần đã cho:
\ketqua{S=S_1+S_2=3+5=8 \ (\text{đvdt})}
\begin{hopchuy}
Lưu ý phân biệt: $\displaystyle\int_{-2}^{4}f(x)\dx = S_1-S_2=3-5=-2$, hoàn toàn khác với diện tích $S=8$. Đây là bẫy rất hay gặp trong câu hỏi trắc nghiệm Đúng/Sai.
\end{hopchuy}
\end{loigiai}

\begin{hopvidu}[(Đọc đồ thị hai hàm số)]
Cho hai hàm số $y=f(x)$ và $y=g(x)$ liên tục trên $[0;2]$, có đồ thị như hình vẽ, trên đó $f(x)\ge g(x)$ với mọi $x\in[0;2]$. Biết $\displaystyle\int_0^2 f(x)\dx=6$ và $\displaystyle\int_0^2 g(x)\dx=2$. Tính diện tích hình phẳng giới hạn bởi hai đồ thị và hai đường thẳng $x=0$, $x=2$.
\end{hopvidu}
\begin{loigiai}
Vì $f(x)\ge g(x)$ trên $[0;2]$ nên không cần xét dấu:
\[
S=\int_0^2\big[f(x)-g(x)\big]\dx = \int_0^2 f(x)\dx - \int_0^2 g(x)\dx.
\]
\ketqua{S=6-2=4 \ (\text{đvdt})}
\end{loigiai}

\begin{hopvidu}[(Xác định công thức đúng từ hình vẽ -- ba nghiệm)]
Đồ thị hàm số $y=f(x)$ (đường bậc ba) cắt trục hoành tại ba điểm có hoành độ $x=-2$, $x=1$, $x=3$ như hình vẽ, trong đó $f(x)\le 0$ trên $[-2;1]$ và $f(x)\ge0$ trên $[1;3]$. Hãy viết công thức đúng để tính diện tích hình phẳng giới hạn bởi đồ thị $y=f(x)$, trục $Ox$ trên đoạn $[-2;3]$ (không giải phương trình $f(x)=0$).
\end{hopvidu}

\begin{center}
\begin{tikzpicture}[scale=0.85, declare function={f(\x)=0.25*(\x+2)*(\x-1)*(\x-3);}]
  \begin{axis}[
      axis lines=middle, xlabel={$x$}, ylabel={$y$},
      xmin=-2.8, xmax=3.6, ymin=-3, ymax=3,
      xtick={-2,1,3}, ytick=\empty, width=8.5cm, height=6.5cm,
      axis line style={-Latex}, samples=100]
    \addplot [domain=-2:1, amber, fill=amber, fill opacity=0.3] {f(x)} \closedcycle;
    \addplot [domain=1:3, teal, fill=teal, fill opacity=0.3] {f(x)} \closedcycle;
    \addplot [domain=-2.6:3.4, navy, thick] {f(x)};
  \end{axis}
\end{tikzpicture}
\end{center}

\begin{loigiai}
Vì $f(x)\le0$ trên $[-2;1]$, phần diện tích tương ứng là $\displaystyle-\int_{-2}^1 f(x)\dx$; vì $f(x)\ge0$ trên $[1;3]$, phần còn lại là $\displaystyle\int_1^3 f(x)\dx$.
\ketqua{S=-\int_{-2}^1 f(x)\dx + \int_1^3 f(x)\dx}
\end{loigiai}

\begin{hopvidu}[(Miền gạch sọc giữa hai đường cong đổi dấu)]
Cho hai đồ thị $y=f(x)$, $y=g(x)$ như hình vẽ (miền gạch sọc), biết trên $[0;2]$ thì $f(x)\ge g(x)$, còn trên $[2;5]$ thì $f(x)\le g(x)$. Viết công thức tính diện tích toàn bộ miền gạch sọc.
\end{hopvidu}

\begin{center}
\begin{tikzpicture}[scale=0.85, declare function={f(\x)=0.3*(\x-2)^2+1; g(\x)=-0.35*(\x-2)^2+3.2;}]
  \begin{axis}[
      axis lines=middle, xlabel={$x$}, ylabel={$y$},
      xmin=-0.5, xmax=5.6, ymin=-0.5, ymax=4,
      xtick={0,2,5}, ytick=\empty, width=8.5cm, height=6cm,
      axis line style={-Latex}, samples=100]
    \addplot [domain=0:2, pattern=north east lines, pattern color=teal] {f(x)} \closedcycle;
    \addplot [domain=0:2, draw=none, fill=white] {g(x)} \closedcycle;
    \addplot [domain=2:5, pattern=north west lines, pattern color=amber] {f(x)} \closedcycle;
    \addplot [domain=2:5, draw=none, fill=white] {g(x)} \closedcycle;
    \addplot [domain=-0.3:5.5, navy, thick] {f(x)} node[right]{$f$};
    \addplot [domain=-0.3:5.5, teal, thick] {g(x)} node[right]{$g$};
  \end{axis}
\end{tikzpicture}
\end{center}

\begin{loigiai}
Trên mỗi khoảng ta lấy (trên trừ dưới) sao cho luôn dương:
\ketqua{S=\int_0^2\big[f(x)-g(x)\big]\dx + \int_2^5\big[g(x)-f(x)\big]\dx}
\end{loigiai}

\newpage
\subsection{Dạng 4: Mô hình hóa toán học (bài toán thực tiễn)}

\begin{hopphuongphap}
\begin{enumerate}[leftmargin=1cm]
  \item Chọn hệ trục tọa độ $Oxy$ hợp lý (thường đặt gốc tại tâm đối xứng hoặc chân công trình) để biểu thức hàm số đơn giản nhất.
  \item Dùng các dữ kiện hình học (chiều cao, chiều rộng, điểm đi qua) để xác định hàm số $y=f(x)$ mô tả biên của công trình.
  \item Tính diện tích bằng tích phân, sau đó chuyển đổi sang đại lượng thực tế cần tìm (chi phí, thể tích đất đào, lượng vật liệu,\ldots).
\end{enumerate}
\end{hopphuongphap}

\begin{hopvidu}[(Cổng chào hình Parabol)]
Một cổng chào có hình dạng Parabol, chiều cao tại đỉnh là $4$ m, chiều rộng chân cổng là $8$ m. Người ta muốn sơn toàn bộ mặt cổng, biết $1$ lít sơn phủ được $5$ m$^2$. Hỏi cần dùng ít nhất bao nhiêu lít sơn?
\end{hopvidu}

\begin{center}
\begin{tikzpicture}[scale=0.75, declare function={f(\x)=-0.25*\x^2+4;}]
  \begin{axis}[
      axis lines=middle, xlabel={$x$ (m)}, ylabel={$y$ (m)},
      xmin=-5.5, xmax=5.5, ymin=-0.7, ymax=4.7,
      xtick={-4,0,4}, ytick={4}, width=8.5cm, height=6cm,
      axis line style={-Latex}, samples=100]
    \addplot [domain=-4:4, navy, fill=navylight, fill opacity=0.35, thick] {f(x)} \closedcycle;
  \end{axis}
\end{tikzpicture}
\end{center}

\begin{loigiai}
Đặt hệ trục tọa độ với gốc $O$ tại trung điểm chân cổng, trục $Oy$ đi qua đỉnh. Đỉnh Parabol là $(0;4)$, chân cổng là hai điểm $(\pm4;0)$ (vì rộng $8$ m). Parabol có dạng $y=ax^2+4$; thay $x=4,y=0$: $0=16a+4 \Rightarrow a=-\dfrac14$.

Vậy $y=-\dfrac14x^2+4$. Diện tích mặt cổng:
\begin{align*}
S&=\int_{-4}^{4}\left(-\frac14x^2+4\right)\dx = 2\int_0^4\left(-\frac14x^2+4\right)\dx \\
 &= 2\left.\left(-\frac{x^3}{12}+4x\right)\right|_0^4 = 2\left(-\frac{16}{3}+16\right)=\frac{64}{3}\text{ m}^2.
\end{align*}
Số lít sơn cần dùng: $\dfrac{64/3}{5}=\dfrac{64}{15}\approx4{,}27$ lít.
\ketqua{\text{Cần ít nhất } 5 \text{ lít sơn (làm tròn lên)}}
\end{loigiai}

\begin{hopvidu}[(Cửa sổ hình bán Elip)]
Một cửa sổ gồm phần thân hình chữ nhật rộng $2{,}4$ m, cao $1{,}5$ m, phía trên gắn thêm một tấm kính hình bán elip có bán trục ngang $a=1{,}2$ m và bán trục đứng (chiều cao vòm) $b=0{,}8$ m. Tính tổng diện tích kính cần lắp cho cả cửa sổ.
\end{hopvidu}
\begin{loigiai}
Đặt gốc tọa độ tại tâm cạnh trên hình chữ nhật, phương trình nửa elip phía trên trục hoành là $y=b\sqrt{1-\dfrac{x^2}{a^2}}$ với $x\in[-a;a]$. Diện tích phần bán elip:
\[
S_{\text{elip}} = \int_{-a}^{a} b\sqrt{1-\frac{x^2}{a^2}}\dx = \frac{\pi ab}{2} = \frac{\pi\cdot1{,}2\cdot0{,}8}{2}=0{,}48\pi\approx1{,}508 \text{ m}^2.
\]
Diện tích phần hình chữ nhật: $S_{\text{cn}}=2{,}4\times1{,}5=3{,}6$ m$^2$.
\ketqua{S_{\text{tổng}} = 3{,}6+0{,}48\pi \approx 5{,}11 \text{ m}^2}
\end{loigiai}

\begin{hopvidu}[(Thể tích đất đào của đường hầm)]
Một đường hầm có mặt cắt ngang dạng vòm Parabol, chiều rộng đáy $10$ m, chiều cao vòm $6$ m, chiều dài hầm $100$ m. Tính khối lượng đất cần đào (theo m$^3$) để tạo đường hầm.
\end{hopvidu}
\begin{loigiai}
Đặt gốc tọa độ tại trung điểm đáy vòm. Đỉnh vòm $(0;6)$, hai chân vòm $(\pm5;0)$. Parabol: $y=ax^2+6$, thay $x=5,y=0$: $a=-\dfrac{6}{25}$.

Diện tích mặt cắt ngang:
\begin{align*}
S&=\int_{-5}^5\left(-\frac{6}{25}x^2+6\right)\dx = 2\int_0^5\left(-\frac{6}{25}x^2+6\right)\dx \\
 &= 2\left(-\frac{6\cdot125}{75}+30\right)=2(-10+30)=40 \text{ m}^2.
\end{align*}
Thể tích đất đào bằng diện tích mặt cắt nhân chiều dài hầm:
\ketqua{V=S\times100=40\times100=4000 \text{ m}^3}
\end{loigiai}

\begin{hopvidu}[(Chi phí lát gạch theo diện tích cong)]
Một hồ cảnh trang trí có viền là hai đường cong $y=\sqrt{x}$ và $y=\dfrac{x}{2}$ (đơn vị mét), với $x\in[0;4]$. Người ta lát gạch phần diện tích giới hạn bởi hai đường cong này với đơn giá $850\,000$ đồng/m$^2$. Tính tổng chi phí lát gạch.
\end{hopvidu}
\begin{loigiai}
Giao điểm: $\sqrt{x}=\dfrac{x}{2} \Leftrightarrow 4x=x^2 \Leftrightarrow x=0$ hoặc $x=4$. Tại $x=1$: $\sqrt1=1 > \dfrac12$, vậy $\sqrt{x}$ nằm trên $\dfrac{x}{2}$ trên $(0;4)$.
\[
S=\int_0^4\left(\sqrt x - \frac x2\right)\dx = \left.\left(\frac23 x^{3/2}-\frac{x^2}{4}\right)\right|_0^4 = \frac{16}{3}-4=\frac43 \text{ m}^2.
\]
Chi phí:
\ketqua{T = \frac43\times850\,000 \approx 1\,133\,333 \text{ đồng}}
\end{loigiai}

\newpage
\subsection{Bài tập tự luyện}

\subsubsection*{Phần I -- Trắc nghiệm nhiều phương án lựa chọn}

\begin{multicols}{2}
\cauhoi Diện tích hình phẳng giới hạn bởi đồ thị $y=x^2-1$, trục $Ox$ và hai đường thẳng $x=0,x=2$ bằng
\begin{enumerate}[label=\textbf{\Alph*.}, leftmargin=0.9cm]
  \item $\dfrac{4}{3}$
  \item $2$
  \item $\dfrac{8}{3}$
  \item $\dfrac{2}{3}$
\end{enumerate}

\cauhoi Diện tích hình phẳng giới hạn bởi $y=\cos x$, trục $Ox$, $x=0$, $x=\pi$ bằng
\begin{enumerate}[label=\textbf{\Alph*.}, leftmargin=0.9cm]
  \item $0$
  \item $1$
  \item $2$
  \item $\pi$
\end{enumerate}

\cauhoi Diện tích hình phẳng giới hạn bởi hai đường $y=x^2$ và $y=4$ bằng
\begin{enumerate}[label=\textbf{\Alph*.}, leftmargin=0.9cm]
  \item $\dfrac{16}{3}$
  \item $\dfrac{32}{3}$
  \item $8$
  \item $16$
\end{enumerate}

\cauhoi Diện tích hình phẳng giới hạn bởi $(P): y=x^2-4x$ và trục hoành bằng
\begin{enumerate}[label=\textbf{\Alph*.}, leftmargin=0.9cm]
  \item $\dfrac{32}{3}$
  \item $\dfrac{16}{3}$
  \item $10$
  \item $\dfrac{64}{3}$
\end{enumerate}

\cauhoi Cho $\displaystyle\int_1^3 f(x)\dx=-5$ và trên $[1;3]$, $f(x)<0$. Diện tích hình phẳng giới hạn bởi $y=f(x)$, $Ox$, $x=1$, $x=3$ bằng
\begin{enumerate}[label=\textbf{\Alph*.}, leftmargin=0.9cm]
  \item $5$
  \item $-5$
  \item $0$
  \item Không xác định được
\end{enumerate}

\cauhoi Diện tích hình phẳng giới hạn bởi $y=e^x$, trục $Ox$, $Oy$ và $x=1$ bằng
\begin{enumerate}[label=\textbf{\Alph*.}, leftmargin=0.9cm]
  \item $e$
  \item $e-1$
  \item $e+1$
  \item $1$
\end{enumerate}

\cauhoi Parabol $y=-x^2+2x$ và trục hoành giới hạn một hình phẳng có diện tích bằng
\begin{enumerate}[label=\textbf{\Alph*.}, leftmargin=0.9cm]
  \item $\dfrac{4}{3}$
  \item $\dfrac{2}{3}$
  \item $2$
  \item $\dfrac{1}{3}$
\end{enumerate}

\cauhoi Diện tích hình phẳng giới hạn bởi $y=x^3$ và $y=x$ (toàn bộ, hai miền) bằng
\begin{enumerate}[label=\textbf{\Alph*.}, leftmargin=0.9cm]
  \item $\dfrac12$
  \item $\dfrac14$
  \item $1$
  \item $\dfrac13$
\end{enumerate}

\cauhoi Cho $(P):y=x^2$ và đường thẳng $d:y=2x$. Diện tích hình phẳng giới hạn bởi $(P)$ và $d$ bằng
\begin{enumerate}[label=\textbf{\Alph*.}, leftmargin=0.9cm]
  \item $\dfrac43$
  \item $\dfrac83$
  \item $4$
  \item $2$
\end{enumerate}

\cauhoi Một cổng Parabol cao $5$ m, rộng chân $6$ m. Diện tích mặt cổng bằng
\begin{enumerate}[label=\textbf{\Alph*.}, leftmargin=0.9cm]
  \item $20$ m$^2$
  \item $15$ m$^2$
  \item $30$ m$^2$
  \item $10$ m$^2$
\end{enumerate}
\end{multicols}

\subsubsection*{Phần II -- Trắc nghiệm Đúng/Sai}

\cauhoi Cho hàm số $y=f(x)=x^2-2x-3$ liên tục trên $\mathbb{R}$, xét trên đoạn $[-2;4]$.
\begin{enumerate}[label=\textbf{\alph*)}, leftmargin=1cm]
  \item Phương trình $f(x)=0$ có hai nghiệm là $x=-1$ và $x=3$.
  \item Trên đoạn $[-1;3]$ thì $f(x)\ge 0$.
  \item Diện tích hình phẳng giới hạn bởi $y=f(x)$, $Ox$, $x=-2$, $x=4$ bằng $\displaystyle\int_{-2}^4 f(x)\dx$.
  \item Diện tích hình phẳng giới hạn bởi $y=f(x)$, $Ox$, $x=-1$, $x=3$ bằng $\dfrac{32}{3}$.
\end{enumerate}

\cauhoi Cho hai hàm số $y=f(x)=\sqrt{x}$ và $y=g(x)=x^2$ xét trên đoạn $[0;1]$.
\begin{enumerate}[label=\textbf{\alph*)}, leftmargin=1cm]
  \item Hai đồ thị cắt nhau tại $x=0$ và $x=1$.
  \item Trên $(0;1)$, $f(x)\ge g(x)$.
  \item Diện tích hình phẳng giới hạn bởi hai đồ thị bằng $\displaystyle\int_0^1(x^2-\sqrt x)\dx$.
  \item Diện tích hình phẳng giới hạn bởi hai đồ thị bằng $\dfrac13$.
\end{enumerate}

\subsubsection*{Phần III -- Trắc nghiệm trả lời ngắn}

\cauhoi Tính diện tích hình phẳng giới hạn bởi đồ thị $y=\ln x$, trục $Ox$ và đường thẳng $x=e$ (làm tròn đến hàng phần trăm).

\cauhoi Một mái vòm nhà kính có mặt cắt hình Parabol cao $3$ m, rộng chân $10$ m. Tính diện tích mặt cắt (m$^2$, làm tròn đến hàng phần chục).

\cauhoi Cho $(P):y=-x^2+4x$ và đường thẳng $d:y=x$. Tính diện tích hình phẳng giới hạn bởi $(P)$ và $d$ (kết quả dạng phân số tối giản).

\cauhoi Một tấm biển quảng cáo hình bán elip có bán trục ngang $2$ m và bán trục đứng $1{,}5$ m được sơn phủ, đơn giá $120\,000$ đồng/m$^2$. Tính chi phí sơn (làm tròn đến nghìn đồng).

\clearpage
\begin{hopchuy}[Đáp án]
\textbf{Phần I:} 1-B, 2-C, 3-B, 4-A, 5-A, 6-B, 7-A, 8-A, 9-A, 10-A.

\textbf{Phần II:} Câu 11: a-Đ, b-S ($f(x)\le0$ trên $[-1;3]$), c-S, d-Đ. \quad Câu 12: a-Đ, b-Đ, c-S (phải lấy $f-g$), d-Đ.

\textbf{Phần III:} Câu 13: $S=1$ (đvdt) $\Rightarrow 1{,}00$. \quad Câu 14: $S=20$ m$^2$. \quad Câu 15: $S=\dfrac{9}{2}$. \quad Câu 16: $S=\dfrac12\pi\cdot2\cdot1{,}5=1{,}5\pi$ m$^2$, $T=1{,}5\pi\times120\,000\approx 565\,000$ đồng.
\end{hopchuy}

\newpage
